\section{Managerial and Operational Interpretation}

The managerial object in this paper is the service menu, not only the final accepted ride. A platform operator controls which meeting-point and pickup-window bundles are visible, how those bundles are priced, and which eligibility filters act before passengers see the menu. The strengthened framework therefore asks managers to audit three layers separately: candidate eligibility, assortment construction, and pricing/choice response.

The exact-vs-greedy diagnostics are useful because they turn a common operational compromise into a measurable tradeoff. If greedy forward selection has a small surrogate-value gap and much lower build time, it is a defensible real-time approximation. If the gap becomes large in certain regimes, the operator can selectively raise the exact-enumeration threshold or reserve exact solving for requests with high expected value.

The robust-filtering comparison is equally important. Hard ETA filters can make the downstream assortment solver appear ineffective by removing feasible meeting-point bundles before they are evaluated. Calibrated and interval-overlap filters provide middle ground: they acknowledge ETA uncertainty without treating every noisy point estimate as a hard exclusion. The quantile error diagnostics (Table~\ref{tab:filter_validity_summary}) show that ETA P95 errors reach approximately 3\,600 seconds, confirming that the filtering signal is coarse but operationally conservative: false-negative pruning concentrates in the far distance band ($>$1\,500\,m), meaning that nearby meeting points---which passengers are most likely to accept---are largely preserved. The no-filter variant remains a diagnostic endpoint, useful for measuring how much menu exposure is being suppressed.

The uptake-regime comparison clarifies when menu optimization should be expected to matter. In very low-uptake settings, observed profit gaps mainly diagnose mechanism behavior because passengers rarely accept non-home bundles. In medium and higher-uptake settings, acceptance, non-home uptake, consumer surplus, and price-boundary rates become informative managerial outcomes. Flat markdown achieves higher net profit than the Lambert-W-inspired rule in the high-uptake regime but reduces acceptance; this is a pricing tradeoff, not a dominance result. A platform should therefore evaluate menu policies in behaviorally live regimes before promoting one pricing rule or filtering rule as operationally superior. The cross-instance results are descriptive two-pair checks; the statistical evidence base for menu-optimization value comes from the six-pair RC benchmark and the uptake-regime stress tests.

The profit decomposition (Table~\ref{tab:profit_decomposition}) reveals that menu-optimization gaps are driven by multiple cost channels simultaneously, not by a single dominant factor. When evaluating whether to deploy a filtering or pricing rule, operators should consider four dimensions jointly: (1) the net-profit direction and magnitude, (2) whether acceptance rates change, (3) whether passenger-facing service metrics such as walk distance, pickup deviation, and in-vehicle time change, and (4) whether the total volume of served requests shifts. The all-request consumer surplus and accepted-user consumer surplus (Appendix~\ref{app:consumer_surplus}) can move in opposite directions when a policy changes the acceptance rate; a profit improvement paired with lower acceptance does not automatically indicate better passenger outcomes.

\paragraph{Defensible TR-E claim.}
The strongest defensible claim of this paper is conditional and mechanism-focused: service-menu optimization---including assortment construction, ETA-informed filtering, and common-offset pricing---produces measurable, split-level paired net-profit differences against dispatch-floor baselines and full-display references under a frozen-predictor, replayed-trace evaluation protocol on the RC benchmark. The direction and rough magnitude of these differences are stable across the tested MNL parameterizations, outside-option levels, and filter configurations, but the absolute profit levels, acceptance rates, and welfare outcomes depend on the demand calibration. The paper does not claim that one policy dominates for all operators; instead, it provides a reproducible framework for evaluating menu policies under operator-specified demand assumptions.
